\documentclass[11pt]{article}
\usepackage[margin=1in]{geometry}
\usepackage{amsmath,amssymb}
\usepackage{booktabs}
\usepackage{hyperref}

\title{PGD Attack: Implementation and Analysis}
\author{Assignment 3 - Part 2}
\date{}

\begin{document}
\maketitle

\section{PGD Implementation}

The PGD (Projected Gradient Descent) attack was implemented in \texttt{adversarial\_attack.py}. PGD is an iterative extension of FGSM that applies multiple smaller steps while projecting back onto the $\epsilon$-ball:

\begin{verbatim}
def pgd_attack(model, data, target, criterion, args):
    x_orig = data.detach()
    x_adv = x_orig.clone()
    for _ in range(num_iter):
        x_adv.requires_grad_(True)
        loss = criterion(model(x_adv), target)
        loss.backward()
        with torch.no_grad():
            x_adv = x_adv + alpha * x_adv.grad.sign()
            perturbation = torch.clamp(x_adv - x_orig, -epsilon, epsilon)
            x_adv = torch.clamp(x_orig + perturbation, 0, 1)
        x_adv = x_adv.detach()
    return x_adv
\end{verbatim}

The PGD defense in \texttt{utils.py} adds adversarial examples to the training batch:
\begin{verbatim}
adv_inputs = pgd_attack(model, inputs, labels, criterion, defense_args)
inputs = torch.cat([inputs, adv_inputs], dim=0)
labels = torch.cat([labels, labels], dim=0)
\end{verbatim}

\section{Question 1: Adversarial Loss vs. Adding Adversarial Examples to Batch}

The two defense approaches are \textbf{different}, though related.

\textbf{Adversarial Loss (FGSM-style):}
\[
\mathcal{L}_{\text{total}} = \mathcal{L}(x, y) + \alpha \cdot \mathcal{L}(x_{\text{adv}}, y)
\]
This computes gradients through both terms simultaneously. The clean and adversarial losses are combined before backpropagation, meaning the gradient update accounts for both objectives in a single step with a weighted combination.

\textbf{Adding Adversarial Examples to Batch (PGD-style):}
The batch is doubled by concatenating clean and adversarial examples, then standard cross-entropy is computed over the augmented batch. Effectively:
\[
\mathcal{L}_{\text{batch}} = \frac{1}{2N}\sum_{i=1}^{N}\left[\mathcal{L}(x_i, y_i) + \mathcal{L}(x_{\text{adv},i}, y_i)\right]
\]

\textbf{When they align:} If $\alpha = 1$ in the adversarial loss formulation, the gradient contributions are mathematically equivalent (both terms weighted equally).

\textbf{When they diverge:}
\begin{itemize}
    \item The adversarial loss allows flexible weighting via $\alpha$, enabling emphasis on clean or adversarial performance.
    \item Batch augmentation treats both equally and changes the effective batch size (which affects batch normalization statistics and learning dynamics).
    \item In FGSM loss, the adversarial example is generated using gradients that are also used for the final update (shared computation graph), whereas in batch augmentation, PGD examples are pre-computed and detached.
\end{itemize}

\section{Question 2: FGSM vs. PGD Tradeoff}

\begin{table}[h]
\centering
\begin{tabular}{lcc}
\toprule
\textbf{Aspect} & \textbf{FGSM} & \textbf{PGD} \\
\midrule
Computation & 1 forward + 1 backward & $k$ forward + $k$ backward \\
Attack strength & Weaker (single step) & Stronger (iterative) \\
Training speed & Fast & Slow ($\sim k\times$ slower) \\
Robustness achieved & Moderate & Higher \\
\bottomrule
\end{tabular}
\end{table}

\textbf{FGSM Advantage:} Computational efficiency. FGSM requires only one gradient computation, making adversarial training significantly faster. This is crucial for large-scale training where PGD's iterative nature becomes prohibitively expensive.

\textbf{PGD Advantage:} Stronger attack and better robustness. PGD finds adversarial examples closer to the worst-case perturbation within the $\epsilon$-ball by iteratively refining the attack. Models trained with PGD defense are typically more robust because they are exposed to stronger adversaries during training, following the principle that defending against a stronger attack generalizes to weaker attacks.

\end{document}
